\section{The graph and difference codes of length 120}

For the residual incidence matrix $H=H_2$, equation (1) gives
$HH^{\mathsf T}=I_{21}$ and $C_{120}=\im H^{\mathsf T}$. Thus the
incidence map is an isometry from the point space to this code. The
constant vector satisfies $H^{\mathsf T}e=j$ and $Hj=40e=e$ over $F$,
so $j\in C_{120}$. The common complement theorem proves
\[
 C_{120}\mathbin{\perp\!\oplus}C_{120}^{\perp}=F^{120},
 \qquad C_{120}^{\perp}=\overline C_{120}.
\]
The $21$-dimensional code and its $99$-dimensional dual have the first
length-$120$ parameters and counts in (4).

The supplementary subcode $S$ in the proof of Proposition~2 is exactly
\[
 S=C_{120}\cap j^\perp=H^{\mathsf T}(e^\perp)
   =\langle(A_{120})_i-(A_{120})_1:1\leq i\leq120\rangle.       \tag{7}
\]
To see the last equality, every row difference has sum zero, while the
rows of $A_{120}$ have the same nonzero sum $77=2$. Their span has
dimension $21$, and adjoining any one row to the difference span gives
that whole space. The difference span therefore has dimension at least
$20$, and containment in the kernel of the nonzero sum functional gives
equality. Under the isometry, that kernel is $e^\perp$.
Since $e^{\mathsf T}e=21=0$, the radical of $e^\perp$ is precisely $Fe$.
Consequently
\[
 \operatorname{Hull}(S)=Fj.                                   \tag{8}
\]
The accepted exact difference-code basis and distribution prove
$S=[120,20,30]_3$ and $S^\perp=[120,100,6]_3$, with respectively
$112$ and $43680$ minimum words.

There is also a direct proof of the source's reflexive-code identity.
On $C_{120}=H^{\mathsf T}F^{21}$ the operator $A_{120}+I$ acts as
$H^{\mathsf T}J_{21}$ on point coordinates; on $C_{120}^\perp=\ker H$
it acts as the identity. Its image is therefore
\[
 \im(A_{120}+I)=Fj+C_{120}^\perp=S^\perp.                     \tag{9}
\]
The two summands have dimensions one and $99$ and meet trivially because
$C_{120}$ is LCD. Thus (9) is an equality of the actual codes, not a match
of parameters.

% The full local group bridge is retained below.
\subsection{Full automorphism groups of the graph and difference codes}

Let $\Gamma$ be the fixed length-$120$ graph, write $A=A_{120}$ and $C=C_{120}$, and let $G$ be the supplied group on its coordinates.

\begin{theorem}\label{L050.120.main}
For these actual coordinate spaces,
\[
 \Aut(C)=\Aut(S)=\Aut(\Gamma)=G,\qquad |G|=80640.
\]
Here $G$ is the group on the literal four coordinate generators of the retained degree-$120$ module certificate. Its previously checked abstract identification is $L_3(4):2^2$. The equality and order in this theorem have the separate elementary certificate below; the abstract identification uses that earlier exact group certificate.
\end{theorem}

\begin{lemma}[Recovering a graph from its code]\label{L050.120.projection}
For this matrix, $E=J-A$ is the orthogonal projector onto $C$. Consequently every coordinate permutation preserving $C$ preserves $A$.
\end{lemma}
\begin{proof}
The strongly regular parameters give $A^2=33I+8A+44J$, hence $A^2=-A-J$ over $\mathbb F_3$. Also $AJ=JA=2J$ and $J^2=0$. Therefore
\[
 E^2=J^2-JA-AJ+A^2=J-A=E,\qquad EA=A.
\]
Every row of $J$ lies in $C$: the sum of the rows of $A$ is $77\mathbf1=2\mathbf1$. Thus $\operatorname{im}E\subseteq C$, and $EA=A$ gives the reverse inclusion (using symmetry to identify row and column spaces). Since $E$ is symmetric and idempotent, its kernel is $C^\perp$ and its image is $C$, so $C\cap C^\perp=0$. This determines the projector uniquely. A coordinate permutation preserving $C$ also preserves the standard dot product and hence $C^\perp$; conjugation must fix $E$. It fixes $J$, so it fixes $A=J-E$.
\end{proof}


\begin{lemma}[Recovering the graph from the difference code]\label{L050.120.support}
The collection of supports of the minimum words of $S$ determines $A$.
\end{lemma}
\begin{proof}
The literal generators of $G$ preserve $A$, hence preserve both $C$ and $S$. Indeed, a permuted row difference is again a difference of two rows. The exact replay verifies a supplied word $w$ lies in the actual $S$ and has weight 30. Its orbit under these generators consists of 112 distinct words. The complete distribution in (4) gives exactly 112 minimum words, so this orbit is all the minimum words of $S$. It is closed under negation and has exactly 56 distinct supports. No orbit size is substituted for an unproved completeness assertion: that assertion uses the retained exhaustive count.

Let $\mathcal B$ be these 56 supports and put
\[
 t_{ij}=|\{B\in\mathcal B:i,j\in B\}|.
\]
The replay constructs every orbit word from $w$, then computes every entry of this integer incidence Gram matrix. It gives $t_{ii}=14$, and, on the 7140 unordered distinct pairs, the complete table is
\begin{center}\begin{tabular}{rrrr}\toprule
$A_{ij}$&$t_{ij}$&Number of pairs&Interpretation\\\midrule
0&4&2520&nonedge\\
1&2&3360&edge\\
1&6&1260&edge\\\bottomrule
\end{tabular}\end{center}
Thus, for $i\ne j$, $A_{ij}=1$ if and only if $t_{ij}\ne4$. Any coordinate permutation preserving $S$ preserves Hamming weight and therefore permutes $\mathcal B$; it preserves every $t_{ij}$ and hence $A$.
\end{proof}

\begin{lemma}[An elementary full graph bound]\label{L050.120.fullness}
The retained four permutations generate all of $\Aut(\Gamma)$.
\end{lemma}
\begin{proof}
First the replay verifies that each is a permutation of all 120 coordinates preserving every adjacency entry. Closing the generated permutations under composition gives exactly 80640 distinct elements; only exact finite permutation operations are used.

For the reverse inequality, apply the common refinement theorem \cite[Theorem~6.4]{U}, using the deterministic choice of the first vertex in the smallest nonsingleton cell (lexicographic ties). 
The following are the actual chosen vertices and sizes of their cells just before individualization; vertices in this table are one-based.
\begin{center}\begin{tabular}{rrrr}\toprule
Step&Chosen vertex&Cell size&Previously fixed vertices\\\midrule
1&1&120&none\\
2&5&42&1\\
3&2&8&1,5\\
4&3&2&1,5,2\\\bottomrule
\end{tabular}\end{center}
After the last refinement all 120 cells are singletons. The common theorem gives $|\Aut(\Gamma)|\le120\cdot42\cdot8\cdot2=80640$.
This matches the subgroup lower bound.
\end{proof}


\begin{proof}[Proof of Theorem~\ref{L050.120.main}]
Every graph automorphism permutes rows of $A$ and therefore preserves $C$ and $S$. Lemmas~\ref{L050.120.projection} and~\ref{L050.120.support} give the reverse inclusions. Lemma~\ref{L050.120.fullness} identifies the common permutation group exactly.
\end{proof}


